\documentclass[11pt,a4paper]{article}
\usepackage[T1]{fontenc}
\usepackage[utf8]{inputenc}
\usepackage[scaled=0.95]{helvet}
\renewcommand{\familydefault}{\sfdefault}
\usepackage[top=1.7cm,bottom=1.5cm,left=2.2cm,right=2.2cm]{geometry}
\usepackage{microtype}
\usepackage[hidelinks]{hyperref}
\pagestyle{empty}
\setlength{\parindent}{0pt}
\setlength{\parskip}{7pt}

\begin{document}

{\large\bfseries Grover Aruquipa}\par\vspace{-6pt}
{\small Qu\'ebec, QC, Canada \ \textbullet\ \ grover.aruquipa.9@gmail.com \ \textbullet\ \ groveraruquipa.github.io}

\vspace{4pt}
Research Team, Robotics and AI Institute, Zurich\\
Research Engineer, Robotics

Dear Research Team,

Most of my time as a PhD student has been spent in the gap between a model that works on paper and a robot that works on the bench. That gap is where your Research Engineers live: connecting simulation, learning, control, software and hardware so that research runs on real robots. It is the work I want to do next, and the Institute is the place where I would most like to do it.

I am finishing my PhD at Universit\'e Laval with Prof.\ Cl\'ement Gosselin and Prof.\ Patrice Lambert, and I defend in early 2027. I designed, built and brought up three multi-DoF robots with 4, 7 and 8 actuators and integrated grippers. For each one I wrote the kinematic and dynamic models, chose the actuators and sensors, wired the motors, drives and encoders, wrote the real-time control in C++ at 500\,Hz, and debugged it on the bench until it met its targets. My 4-DoF coaxial wrist-gripper turned its tool ten full revolutions with less than one degree of tracking error while the gripper held the object.

Simulation is part of how I work, not a separate step. I validate every kinematic model and every trajectory in MuJoCo before it runs on hardware, and I am now training a reinforcement learning policy in NVIDIA Isaac Sim for my Delta robot with an integrated Sarrus gripper. That robot estimates the grasp force from the motor currents, without a sensor in the fingers, and taking its policy from simulation to the real robot is my current project. Your focus on zero-shot sim-to-real deployment on custom hardware is exactly the next step I want to take, with people who do it at the highest level.

My depth is in the parts of the stack that touch the robot: kinematics, dynamics, actuators, sensors, real-time control and the software that ties them together. My breadth comes from the rest of my path. Before my PhD I worked as a computer vision engineer and as a data scientist, including an internship in Zurich with Hitachi ABB Power Grids. I lead two graduate students day to day, so I am used to making technical decisions when a problem is not yet well defined and to working with colleagues who own other parts of the system.

I am available from January 2027 and ready to relocate to Zurich and work on site. Prof.\ Gosselin, Prof.\ Lambert and Univ.-Prof.\ Andreas M\"uller of JKU Linz can speak to my work; their details are in my CV.

Thank you for your time. I would be glad to show you my robots running, on video or live.

Sincerely,\\
Grover Aruquipa

\end{document}
